\documentclass[11pt]{article}
\usepackage[a4paper,margin=25mm]{geometry}
\usepackage{amsmath,amssymb,amsthm}
\usepackage[hidelinks]{hyperref}
\emergencystretch=2em
\newtheorem{theorem}{Theorem}
\title{Karger contraction always finds a minimum cut on a cycle\\
\large Disproof of Conjecture 00000007841}
\author{gaochengzhecpu}\date{}
\begin{document}\maketitle
\begin{abstract}
For every simple cycle $C_n$, $n\geq3$, Karger's
uniform-edge contraction algorithm returns some minimum
cut with probability one. Therefore the claimed bound
$2(1-\varepsilon(n))/n^2$, with $\varepsilon(n)=O(1/n)$,
cannot be sharp on cycles: it tends to zero and its
gap from the actual success probability tends to one.
Lean verifies actual simple input graphs, genuine
multigraph contractions, all histories and normalized
sequential edge-sampling probabilities.
\end{abstract}

\section*{Which success event the source defines}
The source defines $p(G)$ as the probability that the
algorithm outputs \emph{any} minimum cut. It asserts
that its displayed lower bound is tight on cycles;
the Chinese text explicitly claims equality there.
We refute this sharpness clause for arbitrarily large
simple connected cycles. The separate assertion about
the infimum ratio $\rho(n)$ is not assessed here.

Karger's algorithm repeatedly chooses a current
non-loop edge uniformly, merges its endpoints and
deletes loops, retaining parallel edges. At two
supervertices it returns their bipartition, pulled
back to the original vertices. This is the ordinary
contraction algorithm of \cite[Section 2]{karger}.

\section*{An explicit, genuine contraction of a cycle}
Represent the vertices and edge labels of $C_m$ by
$\mathbb Z/m\mathbb Z$, with edge $i$ joining $i$ to
$i+1$. For $m\geq3$ this is the usual simple cycle.
For $m=2$ the same labelled model has two parallel
edges; both must be retained.

Given an edge $e$ of $C_m$, $m\geq3$, rotate it to edge
zero by $x\mapsto x-e$. With representatives
$0\leq x<m$, use the quotient map
\[
 q(x)=\max(x-1,0)\pmod{m-1}.
\]
Only vertices $0$ and $1$ are identified. Edge zero
becomes a loop and is deleted. Each remaining edge
$i=1,\ldots,m-1$ maps to edge $i-1$ of $C_{m-1}$,
with the correct two endpoints. Thus the remaining
edges correspond bijectively, including their
multiplicities, to the next cycle multigraph.

\newpage
\section*{Every contraction history succeeds}
For a Boolean vertex bipartition $A$, its actual
crossing-edge count is
\[
 c_m(A)=\sum_{i\in\mathbb Z/m\mathbb Z}
              \boldsymbol{1}_{A(i)\ne A(i+1)}.
\]
Around a cycle the number of transitions from zero to
one equals the number from one to zero. Consequently
$c_m(A)$ is even. A nonconstant bipartition has a
crossing edge, since equality at every consecutive
pair would make the assignment constant. Every
nontrivial cut therefore has size at least two.

At order two, the final bipartition has two crossing
parallel edges. Pulling a proper cut back through a
surjective contraction preserves properness. It also
preserves its crossing-edge count: the selected edge
has both endpoints in one fiber and contributes zero,
while every other edge corresponds bijectively to a
target edge. Induction backwards through the entire
history shows that every returned cut of the original
cycle is proper and has size two. It is therefore an
actual minimum cut, for every possible history.

At current order $m$ there are exactly $m$ non-loop
edge labels, so every one is selected with probability
$1/m$. The successive choices have orders
$n,n-1,\ldots,3$. The full history set has size
\[
 |\mathcal H_n|=\prod_{m=3}^{n}m=\frac{n!}{2},
\]
and each history has its actual sequential probability
\[
 w(h)=\prod_{m=3}^{n}\frac1m=\frac2{n!}.
\]
These weights are nonnegative and sum to one.
The success event is the assertion that the returned
proper cut is no larger than \emph{every} other proper
cut. Since all histories satisfy this event,
\[
 p(C_n)=\sum_{h\in\mathcal H_n}
        w(h)\boldsymbol{1}_{\text{success}(h)}=1
       \qquad(n\geq3).
\]
This probability is derived from the actual sampling
rule and actual returned partitions; it is not assigned
the value one by definition.

\newpage
\section*{The claimed cycle sharpness is impossible}
If $\varepsilon(n)=O(1/n)$, then $\varepsilon(n)\to0$.
In fact under the weaker condition
$\varepsilon(n)\to0$, the proposed expression satisfies
\[
 b_n:=\frac{2(1-\varepsilon(n))}{n^2}\longrightarrow0,
 \qquad p(C_n)-b_n\longrightarrow1.
\]
It cannot equal the actual success probability on
cycles, even eventually, and it cannot asymptotically
match it. Lean also gives a quantitative contradiction:
for every eventual bound $|\varepsilon(n)|\leq C/n$
and every cutoff, some larger cycle has $b_n<p(C_n)$.

\begin{theorem}
The lower bound stated in Conjecture 00000007841 is
not sharp on simple cycles, whether sharpness means
its asserted exact equality or asymptotic agreement
with the actual success probability.
\end{theorem}
\begin{proof}
All simple cycles have success probability one,
whereas the displayed expression tends to zero.
Their difference tends to one, not zero.
\end{proof}

The familiar contraction estimate concerns survival
of one \emph{prescribed} minimum cut
\cite[Theorem 2.1]{karger}. On a cycle there are many
different minimum cuts, and every returned cut is
minimum. The source defines success as returning any
one of them. This distinction explains the error.
The disproof targets its cycle-sharpness conjunct;
no claim about the actual extremizers or limit of
$\rho(n)$ is needed.

\section*{Bridges to the actual graph and algorithm}
The Lean input is Mathlib's actual simple connected
\texttt{cycleGraph}. Theorems identify its full
undirected edge set with the cyclic edge labels and
prove equality between its genuine crossing-edge count
and the displayed sum. Vertex and edge counts are
checked for every input order at least three.

The contraction map has exactly the selected edge's
two endpoints as its only nonsingleton fiber.
Theorems verify surjectivity, deletion of exactly that
loop and preservation of every other edge's endpoints
and identity. Intermediate order-two states use the
correct two-edge multigraph, not a simplified single
edge. The returned cut is recursively pulled back
through these same quotient maps.

\newpage
\section*{Formal conclusions and reproducibility}
The Lean type \texttt{History k} starts at order $k+2$.
The simple input \texttt{inputGraph n} of order $n+3$
therefore uses \texttt{History (n+1)}. The theorem
\texttt{cycle\_counterexample} bundles the actual
simple connected input, its counts, success probability
one and minimum output cut size two.

The event \texttt{Success} compares an actual proper
output partition with every proper partition.
The function \texttt{pathWeight} multiplies the true
successive $1/m$ factors. Its nonnegativity and total
mass one are proved explicitly. The definition
\texttt{actualSuccessProbability} sums these actual
weights over that event.
Graph-level correctness is proved by
\texttt{actual\_graph\_success}.
\par
The probability-one theorem is named
\texttt{actual\_cycle\_success\_probability}.
The final limit theorems prove that the gap tends to
one and cannot tend to zero, alongside the quantitative
arbitrary-cutoff contradiction.

Lean 4.19.0 and Mathlib revision\\
\texttt{c44e0c8ee63ca166450922a373c7409c5d26b00b}
are pinned. From the submitted \texttt{lean/} directory:
\begin{verbatim}
lake exe cache get
lake build
lake env lean -DwarningAsError=true Main.lean
\end{verbatim}
The complete project is built afresh and checked with
warnings treated as errors. Printed main-theorem
axiom dependencies contain only standard foundational
axioms; no proof placeholders, custom axioms or
external computation oracles are used.

The supplementary \texttt{verify.py} independently
enumerates contractions by original edge identities,
deleting loops and retaining parallel edges. It checks
$C_3$ through $C_9$ exactly, and complete graphs up to
$K_7$, including the nontrivial result $p(K_4)=4/5$.
This is supplementary finite verification; Lean proves
the full cycle family. Run \texttt{python verify.py}
and \texttt{tectonic main.tex}.

The exact source, provenance, build logs, artifact
hashes and PDF checks are included. An agent developed
the proof and supplied a self-review. The parent
separately reviewed the complete proof and final
artifacts. No external independent review or novelty
is claimed for these standard contraction facts.

\begin{thebibliography}{9}
\bibitem{source} The Last Math Competition, Conjecture
00000007841. Exact bilingual source and provenance included.
\bibitem{karger} D. R. Karger and C. Stein,
\emph{An $\widetilde O(n^2)$ Algorithm for Minimum Cuts},
Section 2 and Theorem 2.1.
\href{https://people.csail.mit.edu/karger/Papers/fastcut.pdf}{Author-hosted paper}.
The reference supplies the standard algorithm and
prescribed-cut context; the complete cycle argument
used in this disproof is given above.
\end{thebibliography}
\end{document}
